\PassOptionsToPackage{dvipsnames,table}{xcolor}
\documentclass[11pt,a4paper]{article}

\usepackage{Act}

\begin{document}
\input{\detokenize{/home/fenarius/Travail/Cours/cpge-info/latex/Macros.tex}}

% Entête de la fiche
\Fiche{Erreurs fréquentes en Python}{\sc itc}

\pythonmode

\begin{tcolorbox}[left=0cm,title=\bf{\faSearch \; Méthode : Comprendre un message d'erreur (Traceback)},colbacktitle=cfond]
Lorsqu'une exception survient, Python interrompt l'exécution et affiche une trace d'appels (\textit{traceback}) qui se lit \textit{de bas en haut} :
\begin{enumerate}
	\item[\ding{182}] Lire la \textit{dernière ligne} en \textit{premier} : elle donne le \textit{le type de l'erreur} et un message explicatif précis (ex : \mintinline{python}{IndexError: list index out of range}).
	\item[\ding{183}] \textit{Consulter la ligne incriminée }.\\ {\small \danger} Dans le cas particulier d'une erreur de syntaxe {\tt SyntaxError}, l'erreur réelle peut se situer à la ligne \textit{précédente} (par exemple si une parenthèse n'a pas été refermée)
\end{enumerate}
\end{tcolorbox}

\begin{tcolorbox}[left=0cm,title=\bf{\faBug \; 1. Erreurs de syntaxe (SyntaxError)},colbacktitle=cfond]
Une erreur de syntaxe empêche l'exécution du script avant même son démarrage, ces erreurs sont parfois repérées par l'éditeur de texte.
\begin{itemize}
	\item[\textbullet] \textit{Oubli des deux-points {\tt :}} à la fin d'un en-tête d'instruction conditionnelle, de boucle ou de fonction :
	\setlength{\multicolsep}{0pt}
	\begin{multicols}{2}
		\textcolor{BrickRed}{\faTimesCircle \; Code erroné :}
		\begin{python}
if x > 0
    print("positif")
def carre(n)
    return n * n
		\end{python}
		\textcolor{OliveGreen}{\faCheckCircle \; Code corrigé :}
		\begin{python}
if x > 0:
    print("positif")
def carre(n):
    return n * n
		\end{python}
	\end{multicols}

	\item[\textbullet] \textit{Confusion entre affectation {\tt =} et comparaison {\tt ==}} :
	\begin{multicols}{2}
		\textcolor{BrickRed}{\faTimesCircle \; Code erroné :}
		\begin{python}
if a = 3:
    print("égal à 3")
		\end{python}
		\textcolor{OliveGreen}{\faCheckCircle \; Code corrigé :}
		\begin{python}
if a == 3:
    print("égal à 3")
		\end{python}
	\end{multicols}

	\item[\textbullet] \textit{Parenthèse, crochet ou guillemet non fermé} :
	\begin{multicols}{2}
		\textcolor{BrickRed}{\faTimesCircle \; Code erroné :}
		\begin{python}
val = (2 * x + 1%)
print(val) # SyntaxError ici !
		\end{python}
		\textcolor{OliveGreen}{\faCheckCircle \; Code corrigé :}
		\begin{python}
val = (2 * x + 1)
print(val)
		\end{python}
	\end{multicols}

	\item[\textbullet] \textit{Noms de variables invalides} : un nom de variable ne peut ni commencer par un chiffre (ex : {\tt 1ere\_valeur}), ni contenir de tiret (ex : {\tt total-points} confondu avec une soustraction), ni être un mot réservé Python ({\tt for}, {\tt in}, {\tt def}, {\tt class}, {\tt return}, {\tt pass}, {\tt while}, {\tt import}, \dots).
\end{itemize}
\end{tcolorbox}

\begin{tcolorbox}[left=0cm,title=\bf{\faBug \; 2. Erreurs d'indentation (IndentationError)},colbacktitle=cfond]
En Python, l'indentation structure les blocs de code. Toute incohérence bloque l'exécution. \\
{\small \important} Toujours utiliser la touche \keys{\tab} pour ajouter des indentations (et la combinaison \keys{\shift} + \keys{\tab} pour en retirer). Na \textit{jamais} tenter d'aligner les blocs en utilisant la barre d'espace.
\end{tcolorbox}

\begin{tcolorbox}[left=0cm,title=\bf{\faBug \; 3. Erreurs de type (TypeError)},colbacktitle=cfond]
Survient lors de l'application d'un opérateur ou d'une fonction à un objet d'un type non approprié. Cette erreur survient notamment dans les cas suivants (liste non exhaustive)
\begin{itemize}
	\item[\textbullet] \textit{opération interdite entre types incompatibles} :
	\setlength{\multicolsep}{0pt}
	\begin{multicols}{2}
		\textcolor{BrickRed}{\faTimesCircle \; Code erroné :}
		\begin{python}
s = "Total : " + 42
# Concaténation str + int
L = [1, 2, 3] + 2
# Addition d'une liste et d'un entier
		\end{python}
		\textcolor{OliveGreen}{\faCheckCircle \; Code corrigé :}
		\begin{python}
s = "Total : " + str(42)
# Conversion explicite
L.append(2)
# Méthode appropriée
		\end{python}
	\end{multicols}

	\item[\textbullet] \textit{Flottant utilisé comme indice ou paramètre de \mintinline{python}{range}} :
	\begin{multicols}{2}
		\textcolor{BrickRed}{\faTimesCircle \; Code erroné :}
		\begin{python}
milieu = len(L) / 2
val = L[milieu] # float !
		\end{python}
		\textcolor{OliveGreen}{\faCheckCircle \; Code corrigé :}
		\begin{python}
milieu = len(L) // 2
val = L[milieu] # int !
		\end{python}
	\end{multicols}

	\item[\textbullet] \textit{variable de type {\tt None}}. L'utilisation d'une \textit{méthode} sur un objet de Python (comme un \textit{append}) ne renvois {\sc pas} l'objet mais {\tt None}, de même une fonction qui ne renvoie pas de résultat, renvoie {\tt None}.
	\begin{multicols}{2}
		\textcolor{BrickRed}{\faTimesCircle \; Code erroné :}
		\begin{python}
L = [3, 1, 2]
L = L.append(4)
# L vaut maintenant None !
print(len(L)) # TypeError !
		\end{python}
		\textcolor{OliveGreen}{\faCheckCircle \; Code corrigé :}
		\begin{python}
L = [3, 1, 2]
L.append(4)
# append modifie L en place
print(len(L)) # Affiche 4
		\end{python}
	\end{multicols}
	{\small \danger} Même piège avec \mintinline{python}{L.sort()} (modifie en place et renvoie \mintinline{python}{None}) à distinguer de \mintinline{python}{sorted(L)} (renvoie une nouvelle liste triée).

	\item[\textbullet] \textbf{\mintinline{python}{object is not callable}} : on utilise la syntaxe d'appel à une fonction {\tt <nom\_fonction>(arguments)} alors que {\tt <nom\_fonction>} n'est pas une fonction. L'exemple typique est d'utiliser {\tt ma\_liste(5)} au lieu de {\tt ma\_liste[5]} 
	\item[\textbullet] \textbf{\mintinline{python}{takes X positional arguments but Y were given}} : mauvais nombre d'arguments passés lors de l'appel d'une fonction.

	\item[\textbullet] \textbf{\mintinline{python}{'tuple' / 'str' object does not support item assignment}} : tentative de modification d'un élément d'un type \textit{immuable} (ex : \mintinline{python}{texte[0] = 'A'} ou \mintinline{python}{mon_tuple[1] = 10}).
\end{itemize}
\end{tcolorbox}

\begin{tcolorbox}[left=0cm,title=\bf{\faBug \; 4. Erreurs d'indice (IndexError)},colbacktitle=cfond]
L'exception \mintinline{python}{IndexError: list index out of range} survient lorsqu'on tente d'accéder à un élément avec un indice $i \notin \intN{-\text{len}(L)}{\text{len}(L)-1}$. Voici quelques exemples :
\begin{itemize}
	\item[\textbullet] \textit{Boucle allant un cran trop loin} :
	\setlength{\multicolsep}{0pt}
	\begin{multicols}{2}
		\textcolor{BrickRed}{\faTimesCircle \; Code erroné :}
		\begin{python}
# Indice len(L) inexistant !
for i in range(len(L) + 1):
    print(L[i])
		\end{python}
		\textcolor{OliveGreen}{\faCheckCircle \; Code corrigé :}
		\begin{python}
# Indices de 0 à len(L)-1
for i in range(len(L)):
    print(L[i])
		\end{python}
	\end{multicols}

	\item[\textbullet] \textit{Accès au successeur \mintinline{python}{L[i+1]}} : la borne supérieure de la boucle doit être réduite d'une unité :
	\begin{multicols}{2}
		\textcolor{BrickRed}{\faTimesCircle \; Code erroné :}
		\begin{python}
# Pour i = len(L)-1, i+1 déborde !
for i in range(len(L)):
    if L[i] == L[i+1]:
        print("doublon")
		\end{python}
		\textcolor{OliveGreen}{\faCheckCircle \; Code corrigé :}
		\begin{python}
# Arrêt à len(L)-2
for i in range(len(L) - 1):
    if L[i] == L[i+1]:
        print("doublon")
		\end{python}
	\end{multicols}

	\item[\textbullet] \textit{Affectation directe sur une liste vide / non dimensionnée} :
	\begin{multicols}{2}
		\textcolor{BrickRed}{\faTimesCircle \; Code erroné :}
		\begin{python}
carres = []
for i in range(10):
    carres[i] = i**2 # IndexError !
		\end{python}
		\textcolor{OliveGreen}{\faCheckCircle \; Code corrigé :}
		\begin{python}
carres = []
for i in range(10):
    carres.append(i**2)
		\end{python}
	\end{multicols}
	{\small \important} Pour pré-allouer une liste de taille $n$, écrire \mintinline{python}{carres = [0] * 10}, puis \mintinline{python}{carres[i] = i**2}.
	
\end{itemize}
Les tranches \mintinline{python}{L[debut:fin]} ne lèvent \textit{jamais} d'\mintinline{python}{IndexError}, même si \mintinline{python}{fin} dépasse la longueur de la liste (l'extraction s'arrête simplement au dernier élément).
\end{tcolorbox}

\begin{tcolorbox}[left=0cm,title=\bf{\faBug \; 5. Erreurs de nom et portée (NameError, UnboundLocalError)},colbacktitle=cfond]
Survient lorsqu'on utilise un identificateur non défini ou non accessible dans la portée courante.
\begin{itemize}
	\item[\textbullet] \textbf{\mintinline{python}{NameError: name '...' is not defined}} :
	\begin{itemize}
		\item \textit{Faute de frappe ou casse non respectée :} Python est sensible à la casse (\mintinline{python}{total}, \mintinline{python}{Total} et \mintinline{python}{TOTAL} sont 3 variables distinctes).
		\item \textit{Utilisation avant affectation :} variable manipulée avant sa ligne d'initialisation.
		\item \textit{Oubli d'un import :} utiliser \mintinline{python}{sqrt(x)} sans \mintinline{python}{from math import sqrt} préalable.
		\item \textit{Variable locale inaccessible hors de sa fonction :}
		\setlength{\multicolsep}{0pt}
		\begin{multicols}{2}
			\textcolor{BrickRed}{\faTimesCircle \; Code erroné :}
			\begin{python}
def moyenne(a, b):
    res = (a + b) / 2
    return res

moyenne(12, 16)
print(res) # NameError !
			\end{python}
			\textcolor{OliveGreen}{\faCheckCircle \; Code corrigé :}
			\begin{python}
def moyenne(a, b):
    res = (a + b) / 2
    return res

m = moyenne(12, 16)
print(m) # Récupérer le résultat
			\end{python}
		\end{multicols}
		\item \textit{Variable définie dans un bloc conditionnel non exécuté :} si une variable n'est affectée que dans un \mintinline{python}{if}, elle n'existera pas si la condition est fausse (lui donner une valeur par défaut avant le test).
	\end{itemize}

	\item[\textbullet] \textbf{\mintinline{python}{ModuleNotFoundError: No module named '...'}} : nom de bibliothèque mal orthographié (ex : \mintinline{python}{import maths} au lieu de \mintinline{python}{math}).
\end{itemize}
\end{tcolorbox}

\begin{tcolorbox}[left=0cm,title=\bf{\faBug \; 6. Valeurs, arithmétique et assertions (ValueError, ZeroDivisionError)},colbacktitle=cfond]
L'argument est du bon type, mais sa valeur est invalide ou non acceptée pour l'opération demandée.
\begin{itemize}
	\item[\textbullet] \textbf{\mintinline{python}{ZeroDivisionError: division by zero}} : division \mintinline{python}{/}, quotient entier \mintinline{python}{//} ou modulo \mintinline{python}{%} avec un dénominateur nul (cas d'une liste vide lors du calcul d'une moyenne, ou d'un pivot nul).
	
	\item[\textbullet] \textbf{\mintinline{python}{ValueError: invalid literal for int()}} : tentative de conversion d'une chaîne ne représentant pas un entier :
	\setlength{\multicolsep}{0pt}
	\begin{multicols}{2}
		\textcolor{BrickRed}{\faTimesCircle \; Code erroné :}
		\begin{python}
n = int("3.14") # ValueError !
m = int("dix")  # ValueError !
		\end{python}
		\textcolor{OliveGreen}{\faCheckCircle \; Code corrigé :}
		\begin{python}
n = int(float("3.14")) # Vaut 3
m = 10
		\end{python}
	\end{multicols}

	\item[\textbullet] \textbf{\mintinline{python}{ValueError: math domain error}} : argument hors du domaine de validité mathématique d'une fonction (ex : \mintinline{python}{math.sqrt(-4)}, \mintinline{python}{math.log(0)}, \mintinline{python}{math.asin(2)}).

	\item[\textbullet] \textbf{\mintinline{python}{ValueError: not enough / too many values to unpack}} : déballage d'un tuple ou d'une liste avec un nombre incompatible de variables à gauche de l'affectation :
	\begin{multicols}{2}
		\textcolor{BrickRed}{\faTimesCircle \; Code erroné :}
		\begin{python}
x, y = (1, 2, 3) # Trop de valeurs
a, b, c = [10, 20] # Pas assez
		\end{python}
		\textcolor{OliveGreen}{\faCheckCircle \; Code corrigé :}
		\begin{python}
x, y, z = (1, 2, 3)
a, b = [10, 20]
		\end{python}
	\end{multicols}

	\item[\textbullet] \textbf{\mintinline{python}{AssertionError}} : levé lorsqu'une instruction \mintinline{python}{assert <condition>} échoue (la précondition testée n'est pas vérifiée lors de l'appel).
\end{itemize}
\end{tcolorbox}

\begin{tcolorbox}[left=0cm,title=\bf{\faBug \; 7. Fichiers, dictionnaires et attributs (KeyError, AttributeError)},colbacktitle=cfond]
\begin{itemize}
	\item[\textbullet] \textbf{\mintinline{python}{FileNotFoundError: [Errno 2] No such file or directory}} : tentative d'ouverture avec \mintinline{python}{open("fichier.txt")} d'un fichier inexistant, mal nommé (ex : faute dans l'extension) ou situé dans un autre répertoire (problème de chemin relatif).

	\item[\textbullet] \textbf{\mintinline{python}{KeyError: 'cle'}} : accès direct à une clé qui n'existe pas dans le dictionnaire :
	\setlength{\multicolsep}{0pt}
	\begin{multicols}{2}
		\textcolor{BrickRed}{\faTimesCircle \; Code erroné :}
		\begin{python}
d = {"a": 1, "b": 2}
val = d["c"] # KeyError !
		\end{python}
		\textcolor{OliveGreen}{\faCheckCircle \; Code corrigé :}
		\begin{python}
d = {"a": 1, "b": 2, "c" : 3}
val = d["c"] # "c" est bien une clé !
		\end{python}
	\end{multicols}

	\item[\textbullet] \textbf{\mintinline{python}{AttributeError: '...' object has no attribute '...'}} :\\
	Utilisation d'une méthode inexistante sur le type manipulé (ex : \mintinline{python}{L.apend(5)} au lieu de \mintinline{python}{L.append(5)}).
\end{itemize}
\end{tcolorbox}

\begin{tcolorbox}[left=0cm,title=\bf{\faBug \; 8. Récursion et boucles infinies (RecursionError et terminaison)},colbacktitle=cfond]
\begin{itemize}
	\item[\textbullet] \textbf{\mintinline{python}{RecursionError: maximum recursion depth exceeded}} : trop grand nombre d'appels récursifs imbriqués (limite par défaut à 1000) :
	\begin{itemize}
		\item \textit{Oubli du cas de base (condition d'arrêt) :} la fonction s'appelle indéfiniment.
		\item \textit{Pas de récurrence incorrect :} l'argument ne converge pas vers le cas de base (ex : appel \mintinline{python}{f(n)} au lieu de \mintinline{python}{f(n - 1)}).
	\end{itemize}
	\setlength{\multicolsep}{0pt}
	\begin{multicols}{2}
		\textcolor{BrickRed}{\faTimesCircle \; Code erroné :}
		\begin{python}
def fact(n):
    # Oubli du cas de base !
    return n * fact(n - 1)
		\end{python}
		\textcolor{OliveGreen}{\faCheckCircle \; Code corrigé :}
		\begin{python}
def fact(n):
    if n == 0:
        return 1
    return n * fact(n - 1)
		\end{python}
	\end{multicols}

	\item[\textbullet] \textbf{Boucle \mintinline{python}{while} infinie} : le programme s'exécute sans fin et bloque sans message d'erreur (utiliser alors la combinaison de touche \keys{\ctrl}+\keys{c} pour interrompre le programme):
	\begin{itemize}
		\item Oubli de la mise à jour de la variable de contrôle dans le corps du \mintinline{python}{while} (ex : omission de \mintinline{python}{i += 1}).
		\item Condition d'arrêt stricte sur un flottant : \mintinline{python}{while x != 1.0:} avec un pas \mintinline{python}{x += 0.1} peut ne jamais s'arrêter en raison des erreurs d'arrondi de l'arithmétique flottante.
	\end{itemize}
\end{itemize}
\end{tcolorbox}

\end{document}
